\documentclass[11pt,a4paper]{article}
\usepackage[margin=2.2cm]{geometry}
\usepackage[T1]{fontenc}
\usepackage{amsmath,amssymb,xcolor,enumitem,booktabs,graphicx,tabularx}
\usepackage[normalem]{ulem}
\setlength{\parskip}{0.6em}\setlength{\parindent}{0pt}
\definecolor{newc}{HTML}{1f6f43}
\newcommand{\new}[1]{\textcolor{newc}{#1}}
\newcommand{\old}[1]{\textcolor{gray}{\sout{#1}}}
\newcommand{\note}[1]{{\small\itshape [#1]}}
\newcommand{\lab}[1]{\textsf{\footnotesize\detokenize{#1}}}
\renewcommand{\ref}[1]{\lab{#1}}
\renewcommand{\eqref}[1]{(\lab{#1})}
\newcommand{\nref}[1]{\lab{#1}}
\newcommand{\citep}[1]{[\lab{#1}]}
\newcommand{\citet}[1]{\lab{#1}}
\newcommand{\mockup}[1]{\fbox{\small mockup: #1}}
\title{Proposed change, the Implementation chapter}
\date{25 November 2026, version 1}
\begin{document}\sloppy
\maketitle

\section*{How to read this}
This answers what is left of \texttt{implementation\_comments.md}. Each point
was checked again against \texttt{slap-mapd-coupling} on \texttt{main}
(\texttt{365ccc2}), because the comments are from 20 November.

\textbf{Already settled, not in this review.} A4, the load levels in
\lab{tab:evalparams} now match Appendix~E (0.09, 0.18, 0.3, 0.6, [A111]).
A7, B2, D1 and D2 were applied in [A138]. E1, the timing figure
\lab{app:timing} was never written, and the Results already point to
Appendix~L for where the time goes. B1, the single-agent setup of van
Knippenberg et al.\ stays in ``Choosing together'' only, where it already
is.

Black text stays, \new{green} is new and \old{struck} text goes.
\note{Notes in brackets} say why. Nothing below is in the thesis yet.

\section*{Part 1. What was checked}
\begin{itemize}[nosep]
  \item \texttt{evaluation/evaluator.py:663}: $\kappa_T$ is worked out at the
    end, from the assignment and routing times in the log, added together.
    The service times also come from the task times at the end. The log
    itself is filled while the run goes.
  \item \texttt{environment/multi\_agent\_env.py:227--232}: the backlog is a
    list with one entry per timestep. Blocked is checked every timestep but
    kept as a count per task, plus the longest stretch blocked in a row.
  \item \texttt{training/env.py:537}: for a learned controller the routing
    time includes building each agent's input. \texttt{planners/driver.py:42}:
    for Token Passing it is the planning only.
  \item The timing runs (\texttt{results/L\_results/scaling/step\_time\_main})
    run the network on the GPU, one CPU thread, and Token Passing on the CPU,
    one thread, on the same PC. The Results take runtime from these runs.
  \item \texttt{training/env.py}, docstring: under the centralised and
    section-based controllers a timestep is split into turns. Only the agents
    of a turn are passed in it. The core moves everyone once, after the last
    turn. PettingZoo is not used in the code.
  \item \texttt{training/callbacks.py}: during training the no-conflict check
    and the capacity check run every step and raise an error. The task
    stages hold by how the code is built, and replay is a test-suite check.
    So the four checks cover the hard rules, not the storage rules, the
    inputs or the reward.
  \item ``Trace'' is used in the chapter and in Appendix~K but is never said
    to mean anything.
\end{itemize}

\section*{Part 2. The chapter as it would read}

\subsection*{Simulator}
\ldots\ and \nref{eq:storagegreedy} at every storage update.

All agents \old{act}\new{move} at the same time. \new{Under the centralised
and section-based controllers the agents of a timestep choose in turns
(Section~\ref{sec:method:controllers}), and everyone moves after the last
turn.} Observations, rewards and actions are passed as one entry per agent,
which matches \nref{eq:posg} and the multi-agent \old{interfaces}\new{interface}
of RLlib \old{and PettingZoo} \citep{rllibmultiagent\old{,pettingzooparallel}}.
\note{A6. ``At the same time'' is true for the moves, not the choices, and the
reader has just read about turns. PettingZoo is not used.}

\subsection*{What is logged}
\old{Every measure is added up while the run goes, not computed again at the
end.}\new{What each measure needs is logged while the run goes, so nothing is
worked out again from a replay.} Each task records its release, assignment,
pickup and finish times, whether it is a pick, a put-away or a shipment, and
the order and line it belongs to. The backlog \old{and whether each active
task is blocked are}\new{is} logged every timestep, so a growing queue shows in
the \old{trace}\new{log}. \new{Each task counts the timesteps it is blocked, and
its longest stretch blocked in a row.} Appendix~\ref{app:repro} lists what is
logged and which measure each record gives. \note{A2, the measures are worked
out at the end from what is logged. A3, blocked is a count per task. C2, no
``trace''.}

\subsection*{Computational cost}
Runtime $\kappa_T$ is timed in the simulator. \old{The time per decision is
logged separately for \nref{eq:assignmentrule} and for routing, on the same
machine for every controller.}\new{Assignment and routing are timed apart in
the log, and $\kappa_T$ is their sum, as \nref{eq:runtime} defines it. For a
learned controller routing also holds building each agent's input, because
that is part of producing the joint action. For Token Passing it is the
planning. Runtime is taken from separate timing runs on one machine
(Section~\ref{sec:exp:results}).} The learned controllers run their policy on
the GPU. Token Passing runs on the CPU, the hardware it is built for, so each
method is timed on its own hardware. \note{A5 and C1. One name,
``runtime''. Say that the number is the sum, what routing holds, and which
runs it comes from.}

\subsection*{Correctness and reproducibility}
Four checks test the simulator. \old{The trace has no vertex or swap
conflict}\new{No two agents ever stand on one vertex or swap}, every storage
configuration is within capacity, every task is in exactly one stage at every
timestep, and a run replays exactly under the same seeds. A test suite runs
all four on small warehouses. The first two also run at every timestep of
every training episode, and a run stops if one fails. A task's stage is read
from its times, so it can only have one. These checks show that \old{the
simulator implements the model of Chapter~\ref{sec:problem}}\new{the hard rules
of the model hold}. They are listed in full in Appendix~\ref{app:repro}. Token
Passing in this simulator also gives service times close to those published
by \citet{ma2017lifelong_mapf}, $8\%$ shorter on average
(Appendix~\ref{app:repro:ma}). \note{A8. The checks say nothing about the
storage rules, the inputs or the reward. C2, no ``trace''.}

\section*{Part 3. Appendix K, the same facts}

\lab{app:repro}, first paragraph: \old{Every metric is accumulated online, not
recomputed at the horizon.}\new{What each measure needs is logged while the run
goes, so nothing is worked out again from a replay.} \ldots\ The blocked
indicator of \nref{eq:waiting} is \old{logged as one bit per timestep}\new{checked
every timestep} for the agent serving each active task. \old{The bit}\new{It} is
$1$ when the agent chose $\mathrm{wait}$ or was overridden by the
conflict-resolution operator (Section~\ref{sec:method:resolution}), and
\old{summing it per task gives $W_T$}\new{each task counts these timesteps, which
gives $W_T$}. Backlog $N_t$ is logged every timestep, not only at the horizon,
so a diverging queue is visible in the \old{trace}\new{log}. \note{A2, A3, C2,
so the appendix says what the code does.}

\section*{Questions}
\begin{enumerate}[nosep]
  \item C1. The Results (``Costs'' and the caption of \lab{fig:exp:time}) say
    ``decision time''. Use ``runtime'' there too, so the measure has one name
    from the Problem Formulation to the Results?
  \item A6. PettingZoo is cited only for the shape of the interface. Drop the
    citation, as proposed, or keep it with ``the same shape as''?
\end{enumerate}

\end{document}
